% !TEX root = ../../../main.tex
% Week 4, IMG_9226--IMG_9231. Definitions and examples are expanded from
% the handwritten notes; the limit-law proofs already appear in Chapter 1.
\section{Convergent sequences}
\label{sec:analysis-i:convergent-sequences}

A sequence is a list of real numbers indexed by the natural numbers.
Convergence describes what happens to its terms after sufficiently many
of them have been discarded. The first few terms do not decide its limit.

\begin{definition}
  A sequence $(a_n)$ \emph{converges} to $L\in\R$ if
  \[
    \forall\varepsilon>0\ \exists N\in\N\quad
    n\ge N\ \Longrightarrow\ |a_n-L|<\varepsilon.
  \]
  We write $a_n\to L$ or $\lim_{n\to\infty}a_n=L$.
\end{definition}

The index $N$ may depend on $\varepsilon$. Once it has been chosen,
\emph{every} term with index at least $N$ must lie in
$(L-\varepsilon,L+\varepsilon)$, as in
\autoref{fig:sequence-convergence-band}. Finding a few terms near $L$
would not be enough.

\begin{marginfigure}
  \centering
  % !TEX root = ../../main.tex
% Shared physical dimensions for interval and number-line diagrams.
\tikzset{
  interval diagram/.style={
    x=1cm,y=6mm,font=\footnotesize,
    color=black,line width=0.5pt,line cap=round,line join=round,
    >={Stealth[length=4pt,width=3pt]}
  },
  interval endpoint/.style={
    circle,draw=black,line width=0.5pt,
    minimum size=4pt,inner sep=0pt,outer sep=0pt
  },
  interval open/.style={interval endpoint,fill=white},
  interval closed/.style={interval endpoint,fill=black},
  interval label/.style={inner sep=1pt}
}

\begin{tikzpicture}[interval diagram,x=2.5mm,y=19mm]
  \fill[black!7] (4,0.7) rectangle (12,1.3);
  \draw[gray,dashed] (0,0.7) -- (12,0.7);
  \draw[gray,dashed] (0,1.3) -- (12,1.3);
  \draw[gray] (0,1) -- (12,1);
  \draw[gray,dashed] (4,0) -- (4,1.55);
  \draw[->] (0,0) -- (12.7,0) node[right=4pt] {$n$};
  \draw[->] (0,0) -- (0,1.8) node[above=4pt] {$a_n$};
  \node[left=5pt] at (0,0.7) {$L-\varepsilon$};
  \node[left=5pt] at (0,1) {$L$};
  \node[left=5pt] at (0,1.3) {$L+\varepsilon$};
  \node[below=6pt] at (4,0) {$N$};
  \foreach \n in {1,...,12}{
    \pgfmathsetmacro{\term}{1+(-1)^\n/\n}
    \fill (\n,\term) circle[radius=1.4pt];
  }
\end{tikzpicture}

  \caption[An epsilon band for a convergent sequence.]{For $a_n=1+(-1)^n/n$, all terms from $N=4$ onward lie in
    the band $1-0.3<a_n<1+0.3$. A smaller band may require a later $N$.}
  \label{fig:sequence-convergence-band}
\end{marginfigure}

\subsection{Rules for limits}

If $a_n\to L$, $b_n\to M$, and $c\in\R$, then
\begin{align*}
  ca_n&\to cL, & a_n\pm b_n&\to L\pm M,\\
  a_nb_n&\to LM, & \frac{a_n}{b_n}&\to\frac{L}{M}\quad(M\ne0).
\end{align*}
In the quotient rule the denominator is nonzero for all sufficiently
large $n$, because $b_n\to M\ne0$. These are the rules proved in
\autoref{thm:analysis-i:limit-laws}. The scalar rule follows from the
product rule by using the constant sequence $c$.

% IMG_9226 asks for direct proofs. Present the key estimates as a worked
% example, as requested by the author, rather than a classroom task.
\begin{example}[Bounds in the product and quotient rules]
  \label{ex:analysis-i:limit-law-bounds}
  Suppose $a_n\to L$ and $b_n\to M$. Eventually $|a_n-L|<1$, so
  $|a_n|\le A$ with $A=|L|+1$. The product error satisfies
  \[
    |a_nb_n-LM|\le A|b_n-M|+|M||a_n-L|.
  \]
  Given $\varepsilon>0$, choose $N$ so that this bound on $|a_n|$
  holds and the two errors are respectively less than
  $\varepsilon/(2A)$ and $\varepsilon/(2(|M|+1))$ for $n\ge N$.
  The right-hand side is then less than $\varepsilon$, proving
  $a_nb_n\to LM$ directly from the definition.

  If $M\ne0$, eventually $|b_n-M|<|M|/2$, which gives
  $|b_n|\ge |M|-|b_n-M|>|M|/2$. Hence
  \[
    \left|\frac{a_n}{b_n}-\frac{L}{M}\right|
    \le\frac{2}{|M|}|a_n-L|
      +\frac{2|L|}{|M|^2}|b_n-M|.
  \]
  Choose $N$ so that the denominator bound holds and the two errors
  are respectively less than $\varepsilon|M|/4$ and
  $\varepsilon|M|^2/(4(|L|+1))$ for $n\ge N$. The displayed
  bound is then less than $\varepsilon$, proving the quotient rule.
\end{example}

\begin{theorem}[Squeeze principle]
  \label{thm:analysis-i:squeeze}
  Suppose that $a_n\le b_n\le c_n$ for all sufficiently large $n$.
  If $a_n\to L$ and $c_n\to L$, then $b_n\to L$.
\end{theorem}

\begin{proof}
  Subtracting $L$ gives
  \[
    -|a_n-L|\le a_n-L\le b_n-L\le c_n-L\le |c_n-L|.
  \]
  Consequently $|b_n-L|\le |a_n-L|+|c_n-L|$. Given
  $\varepsilon>0$, choose an index after which the inequalities hold and
  both errors on the right are less than $\varepsilon/2$. Then
  $|b_n-L|<\varepsilon$ for every later term, which proves convergence.
\end{proof}
% IMG_9228 has crossed sign changes. The valid signed-error bounds above
% replace them. IMG_9227's sketch interchanges the upper/lower labels.

The middle sequence may oscillate. What matters is that its available
room shrinks to zero, as \autoref{fig:sequence-squeeze} shows.
In particular, if $0\le a_n\le b_n$ and $b_n\to0$, then $a_n\to0$.

\begin{marginfigure}
  \centering
  % !TEX root = ../../main.tex
% Shared physical dimensions for interval and number-line diagrams.
\tikzset{
  interval diagram/.style={
    x=1cm,y=6mm,font=\footnotesize,
    color=black,line width=0.5pt,line cap=round,line join=round,
    >={Stealth[length=4pt,width=3pt]}
  },
  interval endpoint/.style={
    circle,draw=black,line width=0.5pt,
    minimum size=4pt,inner sep=0pt,outer sep=0pt
  },
  interval open/.style={interval endpoint,fill=white},
  interval closed/.style={interval endpoint,fill=black},
  interval label/.style={inner sep=1pt}
}

\begin{tikzpicture}[interval diagram,x=2.6mm,y=19mm]
  \draw[->] (0,0) -- (12.7,0) node[right=4pt] {$n$};
  \draw[->] (0,0) -- (0,2.25) node[above=4pt] {$y$};
  \draw[gray,dashed] (0,1) -- (12,1);
  \node[left=5pt] at (0,1) {$L$};
  \draw[gray,domain=1:12,samples=80] plot (\x,{1+1/\x});
  \draw[gray,domain=1:12,samples=80] plot (\x,{1-1/\x});
  \node at (4.8,1.55) {$c_n$};
  \node at (4.8,0.45) {$a_n$};
  \foreach \n in {1,...,12}{
    \pgfmathsetmacro{\term}{1+(-1)^\n/(2*\n)}
    \fill (\n,\term) circle[radius=1.4pt];
  }
  \fill (4,1.95) circle[radius=1.4pt];
  \node[right=5pt] at (4,1.95) {$b_n$ terms};
\end{tikzpicture}

  \caption[An oscillating sequence squeezed between two bounds.]{The terms $b_n=1+(-1)^n/(2n)$ are trapped between
    $a_n=1-1/n$ and $c_n=1+1/n$. The gray curves guide the eye;
    the sequence terms are discrete. Both bounds tend to $1$.}
  \label{fig:sequence-squeeze}
\end{marginfigure}

\subsection{Divergence to infinity}

\begin{definition}
  We write $a_n\to+\infty$ if
  \[
    \forall M\in\R\ \exists N\in\N\quad n\ge N\ \Longrightarrow\ a_n>M.
  \]
  Similarly, $a_n\to-\infty$ means that for every $M\in\R$, all
  sufficiently late terms satisfy $a_n<M$.
\end{definition}

These sequences do not have a finite real limit. Divergence to
$+\infty$ means eventually remaining above every prescribed level;
being unbounded alone does not imply this.

\subsection{Geometric sequences}

For a positive number $a$, the behaviour of $a^n$ depends on whether
$a$ is below, equal to, or above $1$.
\[
  a^n\longrightarrow
  \begin{cases}
    0,&0<a<1,\\
    1,&a=1,\\
    +\infty,&a>1.
  \end{cases}
\]
The binomial identity provides a useful estimate. For $n\in\N$,
\[
  (u+v)^n=\sum_{k=0}^n\binom nk u^{n-k}v^k,
  \qquad \binom nk=\frac{n!}{k!(n-k)!}.
\]
If $h>0$, all terms in the expansion of $(1+h)^n$ are nonnegative,
so $(1+h)^n\ge1+nh$.

For $0<a<1$, write $1/a=1+h$ with $h>0$. Then
\[
  0<a^n=\frac{1}{(1+h)^n}\le\frac{1}{1+nh}\longrightarrow0.
\]
The squeeze principle gives the limit. If $a>1$, write $a=1+h$.
Now $a^n\ge1+nh$, which eventually exceeds every $M\in\R$ by the
Archimedean property. Thus $a^n\to+\infty$.

\subsection{Taking increasingly high roots}

\begin{proposition}
  \label{thm:analysis-i:root-sequence-limits}
  We have $n^{1/n}\to1$ and, for every fixed $a>0$, $a^{1/n}\to1$.
\end{proposition}

\begin{proof}
  Set $\delta_n=n^{1/n}-1\ge0$. For $n\ge2$, keeping just the
  quadratic term in the binomial expansion gives
  \[
    n=(1+\delta_n)^n\ge\binom n2\delta_n^2,
    \qquad 0\le\delta_n^2\le\frac{2}{n-1}.
  \]
  Given $\varepsilon>0$, the right-hand side is eventually smaller
  than $\varepsilon^2$, so $0\le\delta_n<\varepsilon$.
  Therefore $n^{1/n}\to1$.

  If $a\ge1$, choose $n>a$. Monotonicity of positive $n$th roots gives
  $1\le a^{1/n}\le n^{1/n}$, so the squeeze principle gives
  $a^{1/n}\to1$. If $0<a<1$, apply this result to $1/a>1$ and use
  the quotient rule to obtain
  \[
    a^{1/n}=\frac{1}{(1/a)^{1/n}}\longrightarrow1.
  \]
\end{proof}

\autoref{fig:root-sequences} plots actual terms of two root sequences.
The base in $n^{1/n}$ grows, yet taking the $n$th root brings the terms
back towards $1$.

\begin{figure*}[htbp]
  \centering
  % !TEX root = ../../main.tex
% Shared physical dimensions for interval and number-line diagrams.
\tikzset{
  interval diagram/.style={
    x=1cm,y=6mm,font=\footnotesize,
    color=black,line width=0.5pt,line cap=round,line join=round,
    >={Stealth[length=4pt,width=3pt]}
  },
  interval endpoint/.style={
    circle,draw=black,line width=0.5pt,
    minimum size=4pt,inner sep=0pt,outer sep=0pt
  },
  interval open/.style={interval endpoint,fill=white},
  interval closed/.style={interval endpoint,fill=black},
  interval label/.style={inner sep=1pt}
}

\begin{tikzpicture}[interval diagram,x=5.5mm,y=25mm]
  \draw[->] (0,0) -- (21.2,0) node[right=5pt] {$n$};
  \draw[->] (0,0) -- (0,2.2) node[above=5pt] {$y$};
  \draw[gray,dashed] (0,1) -- (20.5,1);
  \node[left=5pt] at (0,1) {$1$};
  \node[left=5pt] at (0,2) {$2$};
  \foreach \n in {1,10,20}{
    \draw (\n,0) -- +(0,-2pt);
    \node[below=6pt] at (\n,0) {$\n$};
  }
  \foreach \n in {1,...,20}{
    \pgfmathsetmacro{\growingroot}{\n^(1/\n)}
    \pgfmathsetmacro{\fixedroot}{2^(1/\n)}
    \fill (\n,\growingroot) circle[radius=1.6pt];
    \fill[gray] (\n,\fixedroot) circle[radius=1.6pt];
  }
  \fill (4,-0.5) circle[radius=1.6pt];
  \node[right=6pt] at (4,-0.5) {$n^{1/n}$};
  \fill[gray] (12,-0.5) circle[radius=1.6pt];
  \node[right=6pt] at (12,-0.5) {$2^{1/n}$};
\end{tikzpicture}

  \caption[Two root sequences approaching one.]{Terms of $n^{1/n}$ and $2^{1/n}$ on Cartesian axes.
    Both sequences approach the horizontal line $y=1$. Only integer
    indices represent sequence terms.}
  \label{fig:root-sequences}
\end{figure*}

\subsection{Exercises on root limits}

The following exercises combine root limits with geometric growth and
the squeeze principle.

% !TEX root = ../../hw04.tex
% Source: HW4 Intro to Math Analysis I.pdf, page 1, question 1.
\begin{exercise}
  \label{ex:hw04:polynomial-root-limit}
  \solutionlink{hw04:polynomial-root-limit}
  Find the limit
  \[
    \lim_{n\to\infty}(5n^3-2n+100)^{1/n}.
  \]
\end{exercise}

% !TEX root = ../../hw04.tex
% Source: HW4 Intro to Math Analysis I.pdf, page 1, question 2.
\begin{exercise}
  \label{ex:hw04:exponential-root-limit}
  \solutionlink{hw04:exponential-root-limit}
  Find the limit
  \[
    \lim_{n\to\infty}(a^n+n)^{1/n}
  \]
  for $0\le a\le1$ and $a>1$, respectively.
\end{exercise}

